\documentclass[11pt]{article}

\usepackage[a4paper,margin=2.5cm]{geometry}
\usepackage{amsmath,amssymb,amsthm}
\usepackage{booktabs}
\usepackage{array}
\usepackage{enumitem}
\usepackage[hidelinks]{hyperref}
\usepackage{microtype}

\newtheorem{proposition}{Proposition}
\newtheorem{lemma}{Lemma}
\theoremstyle{definition}
\newtheorem{definition}{Definition}
\newtheorem{assumption}{Assumption}
\theoremstyle{remark}
\newtheorem*{remark}{Remark}

\newcommand{\W}{\mathcal{W}}
\newcommand{\Pdead}{\Pi_{\mathrm{dead}}}
\newcommand{\Tdec}{T_{\mathrm{dec}}}
\newcommand{\Tpre}{T_{\mathrm{pre}}}

\title{A Resource Lower Bound for Single-Stream LLM Inference:\\
the Level-A Reference Line}
\author{Louis Chenluo Deng\\
\texttt{cdeng44@wisc.edu}}
\date{2026-10-09}

\begin{document}
\maketitle

\begin{abstract}
We derive a conditional resource lower bound for single-stream LLM inference and evaluate it for
Qwen3-0.6B on an M4 GPU. Measured device rates give calibrated reference lines, not certified hardware
limits. We retain the exploratory I/O formulas and distinguish them from proved conditional results.
\end{abstract}

\section{Goal}
Profiling shows where an inference engine spends time; a resource reference estimates how much work
is necessary. Following roofline~\cite{roofline}, we use the binding resource and remaining slack to
guide optimization. Scope: Qwen3-0.6B, BF16 weights and KV, M4 GPU, batch 1, decode and prefill.
Engine comparisons and navigation experiments are reported separately in \texttt{docs/report.md}.

\section{Setting}\label{sec:setting}
\subsection{Model}
Qwen3-0.6B is a dense decoder-only transformer with GQA, SwiGLU and tied embeddings.
\begin{center}
\begin{tabular}{lll}
\toprule
Symbol & Meaning & Value \\
\midrule
$\ell$ & layers & 28 \\
$D,\ D_f$ & hidden / FFN size & 1024 / 3072 \\
$H,\ H_{kv},\ d$ & query heads / KV heads / head dimension & 16 / 8 / 128 \\
$V$ & vocabulary & 151\,936 \\
$\beta,\ \tau$ & BF16 element / int32 token bytes & 2 / 4 \\
$\Pi$ & matmul weights per layer & 15\,728\,640 \\
$\nu$ & norm weights per layer (including q/k-norm) & 2304 \\
$\W$ & total weight bytes (tied embedding counted once) & 1\,192\,099\,840 \\
$\kappa$ & KV bytes per token, all layers & 114\,688 \\
\bottomrule
\end{tabular}
\end{center}

\begin{samepage}
\subsection{Device}
MacBook Air M4, 10-core GPU, 24 GB unified memory; only the GPU is modelled.
Rates are microbenchmark calibrations. The spec is \texttt{configs/hardware/m4\_mba\_24g.yaml}.
\begin{center}
\begin{tabular}{lll}
\toprule
Symbol & Resource & Value \\
\midrule
$P$ & FMA throughput & $3.68\times10^{12}$ FLOP/s \\
$B$ & DRAM bandwidth & $102.7\times10^{9}$ B/s \\
$X$ & exp throughput & $4.6\times10^{11}$ exp/s \\
$S$ & modelled fast-memory capacity & $12\ \text{MiB}+10\times512\ \text{KiB}=17\ \text{MiB}$ \\
$P/B$ & machine balance & $35.8$ FLOP/B \\
\bottomrule
\end{tabular}
\end{center}

\end{samepage}

\subsection{Workloads}
\begin{description}[nosep,leftmargin=1.5em]
\item[Decode $(n)$.] Read one input token and $n$ cached positions; attend to $n+1$ keys,
append the input token's KV and sample the next token.
\item[Prefill $(L\ge1)$.] Read $L$ prompt tokens; output their full KV and the first generated token.
\item[Request $(L,N\ge1)$.] Prefill produces token 1; $N-1$ decode steps produce the rest.
\end{description}
All logits are computed, but on-device sampling makes them intermediates: outputs are token ids and
new KV. A full-logits contract adds $V\beta$ bytes per sampled position. Zero-output requests are outside this scope.

\section{Level A: admissible implementations}\label{sec:levelA}
\begin{definition}[Level A]
Fix the reference product DAG. An admissible implementation equals its outputs in exact arithmetic
and computes every useful contraction product at least once, with BF16 weights and KV. Summation
order, fusion, tiling, recomputation and softmax formulation are free.
\end{definition}
For prefill, only the last row of the final layer needs Q, attention, O and MLP; all rows retain their
KV projections. Floating-point differences are checked separately against an FP32 reference.

\section{The bound}\label{sec:bound}
\subsection{Assumptions}
\begin{description}[nosep,leftmargin=1.5em]
\item[(R) Rates.] $P,B,X$ are instantaneous rate ceilings. This is \emph{not certified} for the measured
values above; results are labelled calibrated resource references.
\item[(C) Capacity.] At most $S$ bytes of inputs initially reside on the fast side of DRAM.
$S$ is an upper estimate; overestimating it weakens the bound.
\item[(K) KV.] Cached and output KV are explicit BF16 tensors in DRAM.
\item[(E) Softmax.] Each computed attention score requires one exp.
\item[(T) Tokens.] Input ids are read from, and sampled ids written to, DRAM.
\end{description}

\subsection{Statement}
Let $F$ count useful-product FLOPs (FMA $=2$), $E$ count attention scores, $Q_0$ count compulsory
cold-start DRAM bytes, and $R\le Q_0$ count eligible resident inputs (weights, cached KV and prompt
ids, excluding a decode step's own input token and all output writes).
\begin{proposition}[Resource bound]\label{prop:main}
Under (R)--(T), every level-A implementation of workload $w$ satisfies
\begin{equation}\label{eq:main}
T(w)\ge T_A(w):=\max\left(\frac{F(w)}{P},\frac{Q_0(w)-\min(S,R(w))}{B},\frac{E(w)}{X}\right).
\end{equation}
\end{proposition}
\begin{proof}
Every useful product and required exp must be computed. All weights and cached KV must reach fast
memory, and new KV and sampled ids must reach DRAM. Initial residency saves at most $\min(S,R)$ bytes.
Each resource therefore needs at least its required work divided by its rate; take the maximum.
\end{proof}
The default warm-start credit is $S=17$ MiB, applied once globally (about $0.17$ ms at $B$).
The bound need not be attainable; dispatch and synchronization belong to per-plan costs.

\section{Closed forms for Qwen3-0.6B}\label{sec:closed}
\subsection{Decode, context \texorpdfstring{$n$}{n}}
\begin{gather}
F_{\mathrm{dec}}(n)=2\ell\Pi+4\ell Hd(n+1)+2VD,\qquad E_{\mathrm{dec}}(n)=\ell H(n+1),\\
Q_0=\W+2\tau+(n+1)\kappa,\qquad R=\W+\tau+n\kappa.
\end{gather}
Since $\W>S$ and intensity is about 1--2 FLOP/B, DRAM binds at every context:
\begin{equation}
\boxed{\Tdec(n)=\frac{\W+2\tau-S}{B}+\frac{\kappa}{B}(n+1)
\approx11.4340\ \text{ms}+1.1167\ \mu\text{s}\,(n+1).}
\end{equation}
\begin{center}
\begin{tabular}{rrrr}
\toprule
$n$ & $\Tdec$ & compute time & compute slack \\
\midrule
512 & 12.01 ms & 0.36 ms & $34\times$ \\
4096 & 16.01 ms & 0.58 ms & $28\times$ \\
32768 & 48.03 ms & 2.37 ms & $20\times$ \\
\bottomrule
\end{tabular}
\end{center}
At $n=32768$, KV accounts for about 76\% of compulsory DRAM traffic.

\subsection{Prefill, prompt \texorpdfstring{$L$}{L}}
Let $\Pdead=2DHd+3DD_f=13\,631\,488$ count final-layer Q, O and MLP weights.
\begin{align}
F_{\mathrm{pre}}(L)&=2\ell\Pi L-2\Pdead(L-1)
+4dH\left[(\ell-1)\frac{L(L+1)}{2}+L\right]+2VD,\\
E_{\mathrm{pre}}(L)&=H\left[(\ell-1)\frac{L(L+1)}{2}+L\right],\qquad
Q_0-S=\W+\tau(L+1)+L\kappa-S.
\end{align}
DRAM binds for $L\le48$, compute for $L\ge49$. Attention's quadratic FLOP term exceeds the
linear term near $L=7.7$k and contributes 81\% at 32k.
\begin{center}
\begin{tabular}{rrl}
\toprule
$L$ & $\Tpre$ & binding resource \\
\midrule
16 & 11.45 ms & DRAM \\
64 & 15.06 ms & compute \\
256 & 61.4 ms & compute \\
1024 & 269 ms & compute \\
4096 & 1.454 s & compute \\
32768 & 39.87 s & compute \\
\bottomrule
\end{tabular}
\end{center}
Exp does not bind: $(E/X)/(F_{\mathrm{attn}}/P)=P/(4dX)=1/64$.

\subsection{Requests}
For strictly sequential, non-overlapping phases,
\begin{align}
T_A(L,N)&=\Tpre(L)+\sum_{i=1}^{N-1}\Tdec(L+i-1)\\
&=\Tpre(L)+(N-1)\left[\frac{\W+2\tau-S}{B}+\frac{\kappa}{B}\left(L+\frac{N}{2}\right)\right].
\end{align}
Token dependencies alone do not justify this sum: cross-phase prefetching may overlap work.
General request schedules require separate resource and dependency accounting.

\section{Interpretation for navigation}\label{sec:nav}
\begin{itemize}[nosep]
\item \textbf{Bottleneck and slack:} $\sigma_r=T_A/(u_r/\rho_r)$ identifies resources with spare capacity.
\item \textbf{Gap:} $g=T_{\mathrm{measured}}/T_A$ estimates headroom conditional on the rate ceilings.
\item \textbf{Sensitivity:} for the binding resource, $\partial T_A/\partial\rho_r=-T_A/\rho_r$;
otherwise it is zero away from ties. Using datasheet bandwidth (120 GB/s) lowers decode references by 14.4\%.
\end{itemize}
Closeness to a calibrated line does not establish optimality.

\section{Toward dependency-aware I/O bounds}\label{sec:io}
Finite fast memory may force transfers beyond compulsory traffic. We retain the single-operator
results and exploratory extensions below; none changes the final reference or binding resource in
the evaluated M4 workloads. A tighter whole-graph bound remains open.
Here $M$ is fast memory in elements and $e=\beta$ bytes per element.

\subsection{Matrix multiplication}
\begin{lemma}[Smith et al.~\cite{smith}, Cor.~3.8]\label{lem:gemm}
Classical $C_{m\times n}=A_{m\times k}W_{k\times n}$ with fast memory $M$ and no recomputation reads at least
\begin{equation}
\frac{2mnk}{\sqrt{M}}-mn-2M\quad\text{elements.}
\end{equation}
The leading constant is tight.
\end{lemma}
\paragraph{Exploratory multicore extension.}
With $P$ cores and private capacity $M_c$, the calculator uses
\begin{equation}
\widehat Q_{\mathrm{GEMM}}=\max\left(0,\frac{2mnk}{\sqrt{M_c}}-3PM_c\right)\quad\text{elements.}
\end{equation}
This assumes partitioned classical products, no recomputation and per-core transfer accounting.
Cross-core reductions and charges remain unproved (design notes \S6.8); it is not a general level-A bound.
\paragraph{Intensity cap.}
For large GEMMs, $2mnk/(e\cdot2mnk/\sqrt{M})=\sqrt{M}/e\approx1250$ FLOP/B at 12 MiB
of BF16 storage, about $35\times$ the machine balance.

\subsection{Attention, with explicit constants}
\paragraph{Exploratory refinement.}\label{est:attn}
For $E$ score entries, dimension $d$ and Q/K initially in slow memory, the calculator uses
\begin{equation}
\widehat Q_{\mathrm{attn}}=\max\left(0,\frac{Ed^2}{M+4d^2}-M\right)\quad\text{elements.}
\end{equation}
Inspired by Saha and Ye~\cite{sahaye}, its explicit constants are a local, unproved derivation
(design notes \S6.6).
\paragraph{Proposed extensions.}
Count distinct score pairs for rectangular attention ($E=nL$), causal masks, GQA (stack $g$ query
heads sharing KV, $n=gN$), and multiple heads/layers. Some core computes at least $E/P$ scores;
applying the candidate to that core still requires multicore accounting.
\paragraph{Intensity scale and limitation.}
With $F=4Ed$, the leading term gives $4(M+4d^2)/(de)\approx99000$ FLOP/B at 12 MiB,
and 1280 FLOP/B at 32 KiB ($M=d^2$). In the whole graph, Q/K can be produced or recomputed on chip;
the slow-memory-input assumption therefore restricts this candidate to the operator-boundary estimate.

\subsection{Composing operators}\label{sec:compose}
The calculator's \texttt{--bound operator} and \texttt{--bound hierarchy} estimates sum operator terms.
At boundary $b$, let $P$ count fast-memory instances and $S$ be bytes per instance
(all storage inside the boundary~\cite{savage}): $P=10,S=512$ KiB at the core boundary;
$P=1,S=17$ MiB at DRAM.
\begin{equation}\label{eq:operator}
\widehat Q_b=\max\left(Q_0-PS,\ \sum_o\max(\mathrm{pebble}_o,\mathrm{comp}_o)+Q_{\mathrm{other}}-PS\right).
\end{equation}
Here $o$ covers GEMMs (shared-input projections merged) and attentions; $\mathrm{pebble}_o$ is the
candidate in bytes; $\mathrm{comp}_o$ is cold operator traffic, including activation inputs/outputs
beyond the $PS$ that may remain on chip; and $Q_{\mathrm{other}}$ covers inputs used by neither.
Elementwise operations are treated as fused, with their costs omitted. The global warm-start credit
is applied once, not once per operator.
\begin{equation}\label{eq:hier}
\widehat T=\max\left(\frac{F}{P_{\mathrm{peak}}},\max_b\frac{\widehat Q_b}{B_b},\frac{E}{X}\right).
\end{equation}
$P_{\mathrm{peak}}$ denotes FMA throughput ($P$ in Section~\ref{sec:bound}); $B_b$ is boundary bandwidth.
With only compulsory DRAM traffic $Q_0-\min(S,R)$, this reduces to~\eqref{eq:main}.

\paragraph{Status and limitations.}
$\widehat Q_b$ and $\widehat T$ are exploratory estimates, not proved level-A bounds:
\begin{itemize}[nosep]
\item Fusion can remove operator-boundary transfers; upstream recomputation invalidates local slow-memory-input assumptions.
\item Shared outputs may be charged twice; per-operator residency allowances and the global credit may double-credit bytes.
Summation requires independently charged traffic and consistent residency/sharing accounting, not established here.
\item At prefill $L=4096$, estimated traffic is $8.07\times$ compulsory traffic at the core boundary and
$1.91\times$ at DRAM, yet $\widehat T=T_A$ after taking the resource maximum. Equal times do not imply equal traffic.
\end{itemize}

\subsection{Open problem: whole-graph bounds}\label{sec:notproved}\label{sec:wholegraph}
A whole-graph I/O bound must handle dependencies, arbitrary fusion and recomputation.
Hong--Kung partitioning~\cite{hongkung} is a starting point, subject to the chosen graph and
recomputation model. A feasible fused schedule would supply an upper bound on minimum traffic and
help assess tightness; no schedule attaining $T_A$ is currently established.

\section{Future work}\label{sec:future}
\begin{itemize}[nosep]
\item \textbf{Navigation value.} Test whether resource guidance improves agent optimization over profiling alone.
\item \textbf{Whole-graph I/O.} Establish dependency-aware lower bounds and feasible schedules.
\item \textbf{Extensions.} Study equivalent KV representations, quality-constrained approximation~\cite{harisonak},
MoE, batching, other devices and separate reference lines for speculative or multi-token decoding.
\end{itemize}

\begin{thebibliography}{9}
\bibitem{roofline} S.~Williams, A.~Waterman, D.~Patterson. Roofline: an insightful visual performance
  model for multicore architectures. \emph{Communications of the ACM}, 52(4), 2009.
\bibitem{hongkung} J.-W.~Hong, H.~T.~Kung. I/O complexity: the red-blue pebble game. In \emph{Proc.\ STOC},
  1981.
\bibitem{savage} J.~E.~Savage. Extending the Hong--Kung model to memory hierarchies. In \emph{Proc.\
  COCOON}, 1995.
\bibitem{smith} T.~M.~Smith, B.~Lowery, J.~Langou, R.~A.~van de Geijn. A tight I/O lower bound for
  matrix multiplication. arXiv:1702.02017, 2019.
\bibitem{sahaye} B.~Saha, C.~Ye. The I/O complexity of attention, or how optimal is FlashAttention?
  In \emph{Proc.\ ICML}, 2024. arXiv:2402.07443.
\bibitem{harisonak} T.~Haris, K.~Onak. Compression barriers for autoregressive transformers.
  arXiv:2502.15955, 2025.
\end{thebibliography}

\end{document}
